\title{xLSTM: Extended Long Short-Term Memory}

\begin{abstract}
During the 1990s, the development of Long Short-Term Memory (LSTM) centered around the innovations of the constant error carousel and gating mechanisms. These architectural advances have enabled LSTMs to remain relevant over time, underpinning a multitude of breakthroughs within deep learning—most notably, providing the foundational architecture for early Large Language Models (LLMs). Yet, the emergence of the Transformer model, which is grounded in highly parallelizable self-attention, has signaled a paradigm shift, surpassing LSTM-based approaches in large-scale applications. This work revisits a fundamental question: To what extent can LSTMs advance language modeling performance when scaled to billions of parameters, enhanced by the best practices from contemporary LLMs, and redesigned to overcome their conventional drawbacks? To this end, we present exponential gating, accompanied by normalization and stabilization enhancements. Furthermore, we alter the traditional memory configuration, yielding: (i) the sLSTM variant that features scalar memory with corresponding scalar updates and a novel approach to memory mixing, and (ii) the mLSTM version characterized by a matrix memory and covariance-based update rules, allowing for complete parallelization. By embedding these enhanced LSTM forms within residual block structures, we create xLSTM blocks, which are then composed into xLSTM architectures through residual stacking. These advancements—exponential gating and revised memory formats—significantly extend the capabilities of xLSTM, enabling it to achieve performance and scaling results on par with, and often surpassing, current leading Transformer models and State Space Models.\\
Code available at: \url{https://github.com/NX-AI/xlstm}
\end{abstract}

\section{Introduction}
\label{sec:intro}

The Long Short-Term Memory (LSTM) ideas~\citep{Hochreiter:91,Hochreiter:97nips,Hochreiter:97}, i.e., the constant error carousel and gating, were introduced to overcome the vanishing gradient problem of recurrent neural networks~\citep{Hochreiter:91,Hochreiter:00book}:

\begin{align}
\label{eq:lstm_idea}
  c_t \ = \ f_t \ c_{t-1} \ + \ 
          i_t \  z_t \ , \quad  
          h_t \ = \ o_t \ \psi({c_t}) \ . 
\end{align}

The cell state $c_{t-1}$ (indicated in green) undergoes an additive modification known as the constant error carousel, with cell inputs $z_t$ incorporated and regulated via sigmoid gates (depicted in blue). Specifically, the input gate $i_t$ and the forget gate $f_t$ manage this adjustment to the cell state, whereas the output gate $o_t$ determines the memory cell's output, represented by the hidden state $h_t$. Prior to output gating, the cell state is compressed or normalized by the function $\psi$, and this processed state is subsequently filtered by the output gate to yield the hidden state.

LSTMs have achieved significant success across a wide range of applications~\citep{Hochreiter:01,Hochreiter:07,Schmidhuber:15} and remained the leading architecture for text generation prior to the introduction of Transformers in 2017~\citep{Vaswani:17}. Their utility has been thoroughly established on multiple sequence-based problems, such as text synthesis~\citep{Graves:13,Karpathy:15blog}, handwriting production~\citep{Graves:13}, sequence-to-sequence machine translation~\citep{Sutskever:14nips}, program analysis~\citep{Zaremba:14arxiv}, image captioning~\citep{Karpathy:15,Hossain:19}, automatic code generation~\citep{Karpathy:15blog}, hydrological forecasting such as rainfall-runoff simulation~\citep{Kratzert:18,Kratzert:19}, and flood prediction modeling~\citep{Nearing:24}. In the area of reinforcement learning, LSTMs have served as the top-performing models for sequential data, notably being central components in systems like AlphaStar for StarCraft II~\citep{Vinyals:17short}, OpenAI Five for Dota 2~\citep{OpenAI:19}, and neural controllers in nuclear fusion reactors~\citep{Degrave:22short}. LSTMs are particularly proficient at learning abstract representations—effectively capturing and preserving semantic content within their memory units~\citep{Karpathy:15blog}—as demonstrated, for example, by the discovery of number and syntax neurons~\citep{Lakretz:19}, as well as neurons specializing in linguistic properties~\citep{Bau:19} and sentiment~\citep{Radford:17arxiv}. Despite the emergence of newer models, LSTMs continue to be deployed in pivotal real-world systems~\citep{Degrave:22short,Nearing:24}, underscoring their enduring relevance.

Although LSTMs have achieved considerable progress, they encounter three significant shortcomings: (i) Irrevocable storage decisions. This issue is illustrated with the {\it Nearest Neighbor Search} task (for further details, refer to Appendix~\ref{sec:appNearestNeighborSearch}). In this problem, one seeks to traverse a sequence to locate the vector closest to a reference, in order to return its associated value at the end. As depicted in the left panel of Figure~\ref{fig:lstmProblems}, LSTMs exhibit difficulty updating a previously retained value even after encountering a more similar vector. In contrast, our proposed xLSTM mitigates this shortcoming through exponential gating. (ii) Constrained storage capability, as all information must be compressed into the scalar cell state. This constraint becomes apparent in the {\it Rare Token Prediction} scenario: the right panel of Figure~\ref{fig:lstmProblems} presents the token prediction perplexity on Wikitext-103~\citep{Merity:17} across partitions grouped by token frequency. LSTM underperforms on infrequent tokens due to these capacity constraints, whereas xLSTM, benefiting from matrix-based memory, overcomes this issue. (iii) Inherent lack of parallel processing, a direct consequence of memory mixing via the hidden-to-hidden links across time steps, which necessitates strictly sequential computation. 

Such drawbacks inherent in LSTM architectures have fostered the rise of Transformers~\citep{Vaswani:17} within language modeling. A natural question arises: What new language modeling performance levels can be realized by addressing these challenges and scaling LSTM variants to match the parameter counts of leading Large Language Models?

\section{Extended Long Short-Term Memory}
\label{sec:xLSTM}

To address the shortcomings of LSTM, the Extended Long Short-Term Memory (xLSTM) framework incorporates two significant innovations into the LSTM framework described in Equation~\eqref{eq:lstm_idea}. These innovations—namely exponential gating and the introduction of new memory designs—expand the LSTM model family with two variants: (i) sLSTM (see Section~\ref{sec:sLSTM}), which is characterized by a scalar memory, scalar updates, and memory mixing, and (ii) mLSTM (see Section~\ref{sec:mLSTM}), which employs a matrix-based memory and leverages a covariance-based (outer product) update rule, supporting full parallelism. Both sLSTM and mLSTM benefit from the exponential gating mechanism. In order to permit parallel processing, mLSTM omits memory mixing, specifically the recurrent hidden-to-hidden interactions. Additionally, both architectures allow for extensions to multiple memory cells, with sLSTM supporting memory mixing across these cells. The sLSTM model can also utilize multiple heads, enabling memory mixing across cells within each head, while restricting memory mixing between different heads—a mechanism that, combined with exponential gating, yields an alternative form of memory mixing. In contrast, for mLSTM, multi-head and multi-cell structures are effectively interchangeable.

By incorporating these LSTM extensions within residual block structures, one obtains xLSTM blocks (refer to Section~\ref{sec:xLSTMblock}). Building deep models by stacking such xLSTM blocks residually produces xLSTM-based architectures (see Section~\ref{sec:xLSTMarchitecure}). Figure~\ref{fig:xlstm_sketch} offers an overview of the xLSTM architecture and its principal components.

\subsection{Review of the Long Short-Term Memory}
\label{sec:orgLSTM}

The foundational concept behind the original LSTM~\citep{Hochreiter:91,Hochreiter:97nips,Hochreiter:97} centers around a scalar memory cell that serves both as the primary storage and computational element, addressing the vanishing gradient issue~\citep{Hochreiter:91,Hochreiter:00book} via the so-called constant error carousel mechanism (i.e., how the cell state is updated). This memory cell architecture features three essential gating mechanisms: the input gate, the output gate, and the forget gate—the last of which was later introduced by \citet{Gers:00}. The dynamics of the LSTM memory cell at timestep $t$ can be expressed using the following set of update equations:

\begin{align}
c_t \ &= \  f_t \ c_{t-1} \ + \ i_t \ z_t &  & &\text{cell state} \\
h_t  \ &= \ o_t \ \tilde{h}_t \ , & \tilde{h}_t \ &= \ \psi \left( c_t\right)
  &\text{hidden state} \\
z_t \ &= \ \varphi \left( \tilde{z}_t \right) \ , 
  &\tilde{z}_t \ &=  \ \mathbf{w}^\top_{z} \ \mathbf{x}_t \ + \
  r_{z}  h_{t-1} \ + \  b_{z} \ \
  &\text{cell input} \\
i_t \ &= \ \sigma \left( \tilde{i}_t  \right) \ , 
  &\tilde{i}_t \ &= \ \mathbf{w}^\top_{i} \ \mathbf{x}_t \ + \
  r_{i}  \ h_{t-1} \ + \  b_{i} \ \
  &\text{input gate} \\
f_t \ &= \ \sigma \left(  \tilde{f}_t \right) \ , 
  &\tilde{f}_t \ &= \ \mathbf{w}^\top_{f} \ \mathbf{x}_t  \ + \
  r_{f}  \ h_{t-1} \ + \  b_{f} \ \
  &\text{forget gate} \\
o_t \ &= \ \sigma \left( \tilde{o}_t \right) \ , 
  &\tilde{o}_t  \ &= \ \mathbf{w}^\top_{o} \ \mathbf{x}_t \ + \
  r_{o}  \ h_{t-1} \ + \  b_{o} \ \
  &\text{output gate} 
\end{align}

The weight vectors $\mathbf{w}_{z}$, $\mathbf{w}_{i}$, $\mathbf{w}_{f}$, and $\mathbf{w}_{o}$ denote the connections from the input vector $\mathbf{x}_t$ to the cell input, input gate, forget gate, and output gate, respectively. Recurrent weights $r_{z}$, $r_{i}$, $r_{f}$, and $r_{o}$ link the previous hidden state $h_{t-1}$ to the cell input and each respective gate. The variables $b_{z}$, $b_{i}$, $b_{f}$, and $b_{o}$ represent the biases associated with each component. The cell input and hidden state activations are represented by functions $\varphi$ and $\psi$, often chosen as $\tanh$. The activation function $\psi$ serves to bound the cell state, as without it the values could diverge. All gates employ the sigmoid activation, defined as $\sigma \left( x \right) = 1/(1+\exp(-x))$. Subsequent variants aggregated several scalar memory cells into a single vector $\mathbf{c}_t \in \mathbb{R}^{d}$, thereby enabling the introduction of recurrent weight matrices $\mathbf{R} \in \mathbb{R}^{d \times d}$ that facilitate interaction among different memory cells~\citep{Greff:15}; additional discussion appears in Appendix~\ref{sec:apporgLSTM}. Experiments removing model elements (ablation studies) have demonstrated the necessity of each component within the memory cell~\citep{Greff:15}.

\subsection{sLSTM}
\label{sec:sLSTM}

To enable LSTMs to dynamically adjust their memory allocation, we propose the addition of exponential gates (highlighted in red), as well as incorporate normalization and stabilization techniques. Specifically, we allow the input and forget gates to use exponential activation functions. For the normalization process, we establish a normalizer state that accumulates the sum of the input gate multiplied by the sequence of all subsequent forget gates.

The scalar sLSTM forward pass is:

\begin{align}
\label{eq:slstmforward}
c_t \ &= \  f_t \ c_{t-1} \ + \ i_t \ z_t & & 
 &\text{cell state} \\
n_t \ &= \  f_t \ n_{t-1} \ + \ i_t  & & 
 &\text{normalizer state} \\
h_t \ &= \ o_t \ \tilde{h}_t \ ,
  &\tilde{h}_t \ &= \ c_t / n_t
&\text{hidden state} \\
z_t \ &= \ \varphi \left( \tilde{z}_t \right) \ , 
  &\tilde{z}_t \ &=  \ \mathbf{w}^\top_{z} \ \mathbf{x}_t \ + \
  r_{z}  h_{t-1} \ + \  b_{z}
  &\text{cell input} \\
i_t \ &= \ \!\exp \left( \tilde{i}_t  \right) \ , 
  &\tilde{i}_t \ &= \ \mathbf{w}^\top_{i} \ \mathbf{x}_t \ + \
  r_{i}  \ h_{t-1} \ + \  b_{i}
  &\text{input gate} \\
f_t \ &= \ \sigma \left(  \tilde{f}_t \right) \ \text{OR} \
\!\exp \left(  \tilde{f}_t \right) \ ,
  &\tilde{f}_t \ &= \ \mathbf{w}^\top_{f} \ \mathbf{x}_t  \ + \
  r_{f}  \ h_{t-1} \ + \  b_{f}
  &\text{forget gate} \\
o_t \ &= \ \sigma \left( \tilde{o}_t \right) \ , 
  &\tilde{o}_t  \ &= \ \mathbf{w}^\top_{o} \ \mathbf{x}_t \ + \
  r_{o}  \ h_{t-1} \ + \  b_{o}
  &\text{output gate} 
\end{align}

We transfer the original LSTM gating techniques, i.e., input- and/or hidden-dependent gating plus bias term, to the new architectures. Exponential activation functions can lead to large values that cause overflows. Therefore, we stabilize gates with an additional state \mbox{$m_t$~\citep{Milakov:18arxiv}}:

\begin{align}
\label{eq:slstmstabil}
m_t \ &= \ \max \left( \log ( f_t ) + m_{t-1} , \log ( i_t ) \right) &\text{stabilizer state} \\
i'_t \ &= \ \!\exp \left( \log \left ( i_t \right) - m_t \right) \ = \ \!\exp \left( \tilde{i}_t - m_t \right) \ \, 
  &\text{stabil. input gate} \\
f'_t \ &= \ \!\exp \left( \log \left( f_t \right) + m_{t-1} - m_t \right) \ \, 
  &\text{stabil. forget gate}
\end{align}

Appendix~\ref{sec:appsLSTM} demonstrates that substituting $f_t$ with $f'_t$ and $i_t$ with $i'_t$ during the forward computation does not alter either the network’s overall output or the gradients of the loss relative to its parameters.

\paragraph{New Memory Mixing.} Similar to the conventional LSTM, the sLSTM architecture may incorporate several memory cells (refer to Appendix~\ref{sec:appsLSTM}). Utilizing multiple memory cells facilitates memory mixing, enabled through recurrent links $\mathbf{R}_{\mathbf{z}}$, $\mathbf{R}_{\mathbf{i}}$, $\mathbf{R}_{\mathbf{f}}$, $\mathbf{R}_{\mathbf{o}}$ that connect the hidden state vector $\mathbf{h}$ with the memory cell input $\mathbf{z}$ and the respective gates $\mathbf{i}$, $\mathbf{f}$, and $\mathbf{o}$. A distinctive feature of memory mixing in this context is introduced by the exponential gating mechanism. The redesigned sLSTM supports multiple heads, where memory mixing occurs within individual heads but not between them. The combination of head structures and exponential gating in sLSTM establishes an innovative approach to memory mixing.

\subsection{mLSTM}
\label{sec:mLSTM}

To improve the memory capabilities of LSTMs, we extend the standard memory cell from a single scalar $c \in \mathbb{R}$ to a matrix form $\mathbf{C} \in \mathbb{R}^{d \times d}$. This allows retrieval operations to utilize matrix multiplication for accessing stored information. At a particular time step $t$, our goal is to encode two vectors: the key $\mathbf{k}_t \in \mathbb{R}^d$ and the value $\mathbf{v}_t \in \mathbb{R}^d$—terminology adopted from Transformer models. Subsequently, at time $t + \tau$, we want the value $\mathbf{v}_t$ to be recoverable using a query vector $\mathbf{q}_{t+\tau} \in \mathbb{R}^d$. This architecture corresponds to the framework used in Bidirectional Associative Memories (BAMs) as described in \citep{Kohonen:72,Anderson:72,Nakano:72,Anderson:77}.

The covariance update rule~\citep{Sejnowski:77,Dayan:91} for storing a key-value pair is

\begin{align}
\mathbf{C}_t \ &= \ \mathbf{C}_{t-1} \ + \ \mathbf{v}_{t} \ \mathbf{k}_t^\top \ .
\end{align}

We posit that a layer-norm is applied prior to the projection of inputs into keys and values, ensuring these vectors maintain zero mean. The covariance update mechanism, as described in~\citep{Dayan:91}, achieves optimality in terms of maximizing the separability among retrieved binary vectors; this objective aligns with enhancing the signal-to-noise ratio. When restricting retrieval to only pairwise associations and accepting quadratic computational overhead—as seen in attention mechanisms~\citep{Krotov:16,Krotov:17,Ramsauer:21}—even greater separability can be attained. Furthermore, the covariance update procedure corresponds to the formulation found in Fast Weight Programmers~\citep{Schmidhuber:92ncfastweights,Schlag:21}, later extended through the incorporation of both a fixed decay factor applied to $\mathbf{C}_{t-1}$ and a static learning rate scaled by 
$\mathbf{v}_{t} \mathbf{k}_t ^\top$~\citep{Ba:16}.
Following this conceptual lineage, we adapt the covariance update scheme within the LSTM architecture, mapping the forget gate to a decay parameter, the input gate to a learning parameter, and utilizing the output gate to modulate the scale of the extracted vector.

For this matrix memory, the normalizer state is the weighted sum of key vectors, where each key vector is weighted by the input gate and all future forget gates. Again, the normalizer state keeps record of the strength of the gates. Since the dot product between query and normalizer state can be close to zero, we use the absolute value of this dot product and lower bound it by a threshold (typically 1.0) as done previously~\citep{Sun:23arxiv}. The mLSTM forward pass is:

\begin{align}
\label{eq:mlstm_recurrent_begin}
\mathbf{C}_t \ &= \  f_t \ \mathbf{C}_{t-1} \ + \ 
  i_t \ \mathbf{v}_t \ \mathbf{k}_t^\top & &  &\text{cell state} \\
\mathbf{n}_t \ &= \  f_t \ \mathbf{n}_{t-1} \ + \ 
  i_t \ \mathbf{k}_t & &  &\text{normalizer state} \\
\mathbf{h}_t  \ &= \ \mathbf{o}_t \ \odot \ \tilde{\mathbf{h}}_t
\ , \qquad \qquad 
\tilde{\mathbf{h}}_t \ = \   \mathbf{C}_t \mathbf{q}_t \ / \ 
  \max \left\{ |\mathbf{n}_t^\top \mathbf{q}_t|, 1 \right\} 
   &&&\text{hidden state} \\
\mathbf{q}_t \ &= \ \mathbf{W}_q \ \mathbf{x}_t \ + \ \mathbf{b}_q  & &  &\text{query input} \\
\mathbf{k}_t \ &= \ \frac{1}{\sqrt{d}} \mathbf{W}_k \ \mathbf{x}_t \ + \ \mathbf{b}_k  & &  &\text{key input} \\
\mathbf{v}_t \ &= \ \mathbf{W}_v \ \mathbf{x}_t \ + \ \mathbf{b}_v  & &  &\text{value input} \\
i_t \ &= \ \!\exp \!  \! \left( \tilde{i}_t  \right) \ , 
 \qquad \qquad \ \ \ \ \,
  \tilde{i}_t \ = \ \mathbf{w}^\top_{i} \ \mathbf{x}_t \ + \  b_{i} \ \ \ 
  &&&\text{input gate} \\
f_t \ &= \ \sigma \!  \left(  \tilde{f}_t \right) \ \text{OR} \ \!\exp \!  \! \left(  \tilde{f}_t \right) , \ \ \: \!
\tilde{f}_t \ = \ \mathbf{w}^\top_{f} \ \mathbf{x}_t  \ + \
  b_{f}
  &&& \text{forget gate} \\
\label{eq:mlstm_recurrent_end}
\mathbf{o}_t \ &= \ \sigma \left( \tilde{\mathbf{o}}_t \right) \ , \qquad \qquad \quad \ \ \ \,
  \tilde{\mathbf{o}}_t  \ = \ \mathbf{W}_{\mathbf{o}} \ \mathbf{x}_t \ + \
  \mathbf{b}_{\mathbf{o}}
  &&&\text{output gate} 
\end{align}

Like the traditional LSTM, mLSTM allows for multiple memory cells. In the mLSTM architecture, using several heads or several memory cells is functionally identical due to the absence of interactions between memories. To address potential instability from the exponential gates in mLSTM, we employ stabilization methods identical to those utilized for sLSTM (refer to Equation~\ref{eq:slstmstabil}). Owing to the lack of memory interaction in mLSTM, the recurrence relation can be re-expressed in a parallelized form. Additional information can be found in Appendix~\ref{sec:appmLSTM}.

\subsection{xLSTM Architecture}
\label{sec:xLSTMarch}

\paragraph{xLSTM Blocks.}
\label{sec:xLSTMblock}
The xLSTM block is designed to non-linearly encode previous information into a high-dimensional representation, enhancing the distinction between various histories or contexts. This distinction is crucial for accurate prediction of forthcoming sequence elements such as the next token. Our approach leverages Cover’s Theorem~\citep{Cover:65}, which posits that patterns, when non-linearly projected into a space of higher dimensionality, become more amenable to linear separation compared to the original space. We examine two variants of residual block structures: (i) The first, a residual block featuring a post up-projection (similar to Transformers), first conducts a non-linear aggregation of the historical sequence in its original space, subsequently performs a linear projection to a high-dimensional space, applies a non-linear activation, and finally projects the data linearly back to the original space. This design can be seen in the left section of Figure~\ref{fig:Backbones}
as well as in the third column of Figure~\ref{fig:xlstm_sketch}.
A comprehensive schematic is available in Figure~\ref{fig:appsLSTM_detailed}
in the appendix. (ii) The second variant employs a residual block with a pre up-projection (inspired by State Space Models), wherein a linear map is first performed into a high-dimensional space,

provides a nonlinear abstraction of prior information in a high-dimensional representation, followed by a linear transformation back to the original input space. When an xLSTM block includes an sLSTM, the post up-projection block is employed. Conversely, if the xLSTM block contains an mLSTM, the pre up-projection block is applied, capitalizing on the increased memory capacity afforded by the high-dimensional representation. For further clarification, refer to the left panel of Figure~\ref{fig:Backbones} and the third column of Figure~\ref{fig:xlstm_sketch}, as well as Figure~\ref{fig:appsLSTM_detailed} in the appendix.

\paragraph{xLSTM Architecture. }
\label{sec:xLSTMarchitecure}
The xLSTM architecture is designed through residual stacking of modular units~\citep{Srivastava:15,He:16}. Our implementation follows standard practice in contemporary Large Language Models, utilizing pre-LayerNorm~\citep{Ba:16b} residual architectures. An illustration can be found in the final column of Figure~\ref{fig:xlstm_sketch}.

\subsection{Memory and Speed Considerations}

Unlike Transformers, xLSTM networks maintain both linear computational complexity and constant memory requirements relative to the sequence length. Thanks to the compressive nature of xLSTM’s memory, these models are particularly advantageous for deployment in industrial contexts and on edge devices.

In the case of mLSTM, its memory architecture does not demand additional parameters but incurs significant computational cost, stemming from its $d \times d$ matrix-based memory and corresponding $d \times d$ update operations. This design represents a balance between enhanced memory capability and increased computational expense. However, since these operations can be efficiently parallelized using GPUs, their influence on the overall runtime is minimal.

Although mLSTM can exploit parallel computation, similar to FlashAttention~\citep{Dao:22,Dao:23} and GLA~\citep{Yang:23arxiv}, sLSTM lacks parallelizability due to the intrinsic hidden-hidden memory mixing. Nevertheless, we have engineered an optimized CUDA implementation with memory handling down to the register level, resulting in sLSTM achieving a runtime typically within a factor of two compared to mLSTM.

\section{Related Work}
\label{sec:related}

\paragraph{Linear Attention.} Various approaches have been proposed to address the quadratic dependency of the Transformer’s complexity on the context length, aiming to achieve attention that scales linearly with context. Synthesizer generates synthetic attention scores independently of token--token relationships~\citep{Tay:20arxiv}. Linformer uses a low-rank matrix to implement self-attention and directly provides a linear approximation~\citep{Wang:20arxiv}. Linear Transformer transforms the attention operation into a linear process~\citep{Katharopoulos:20}. Performer approximates the attention softmax function with a method based on positive orthogonal random features~\citep{Choromanski:21}. Fast, extended convolutions replace traditional attention in the Structured Global Convolution (SGConv)~\citep{Li:22} and the Hyena Hierarchy~\citep{Poli:23}.

\paragraph{State Space Models.} In recent years, State Space Models (SSMs) have gained substantial attention, largely due to their linear scalability with respect to context length and their competitive results relative to Transformers. Notable early work in this area includes the introduction of the Structured State Space sequence model (S4)~\citep{Gu:21}. This was subsequently expanded upon by approaches such as the Diagonal State Space (DSS) model~\citep{Gupta:22}, Gated State Space (GSS) models~\citep{Mehta:22}, the S5 model~\citep{Smith:22}, Bidirectional Gated SSM (BiGS)~\citep{Wang:22}, H3 model~\citep{Fu:23}, and Mamba~\citep{Gu:24arxiv}.

\paragraph{Recurrent Neural Networks.} In recent studies, Recurrent Neural Networks (RNNs) have been considered as alternatives to Transformer and attention mechanisms, primarily because their computational complexity scales linearly with respect to the input sequence length. Deep Linear Recurrent Units (LRUs) within RNN architectures have delivered strong performance in language modeling tasks~\citep{Orvieto:23, De:24arxiv}. Hierarchically Gated Linear RNNs (HGRN) \citep{Qin:23} and the subsequent HGRN2~\citep{Qin:24arxiv} have likewise yielded favorable results. Among RNN-based models, RWKV~\citep{Peng:23arxivshort, Peng:24arxivshort} is notable for its efficacy in large language model settings, achieving performance on par with Transformer-based models.

\paragraph{Gating.}
The concept of gating, fundamental to the LSTM architecture, has been revisited and adapted within many recent models. Variants of gating have been integrated in architectures including HGRN~\citep{Qin:23}, HGRN2~\citep{Qin:24arxiv}, Gated Linear Attention (GLA)~\citep{Yang:23arxiv}, Gated State Space (GSS) models~\citep{Mehta:22}, Bidirectional Gated SSM (BiGS)~\citep{Wang:22}, Moving Average Equipped Gated Attention (MEGA)~\citep{Ma:22}, RWKV~\citep{Peng:23arxivshort}, as well as Mamba~\citep{Gu:24arxiv}.

\paragraph{Covariance Update Rule.}
To increase the storage capabilities, we introduced a covariance update mechanism to the mLSTM cell's matrix memory. Several approaches employ similar covariance-based updates, such as Fast Weight Programmers~\citep{Schmidhuber:92ncfastweights,Schlag:21}, RWKV-5 and RWKV-6~\citep{Peng:24arxivshort}, Retention~\citep{Sun:23arxiv}, the Linear Transformer~\citep{Katharopoulos:20}, and HGRN2~\citep{Qin:24arxiv}.

\paragraph{Most Related.} The models that are most analogous to xLSTM in terms of design are Retention~\citep{Sun:23arxiv}, RWKV~\citep{Peng:23arxivshort,Peng:24arxivshort}, and HGRN2~\citep{Qin:24arxiv}, which incorporate matrix memory and/or gating mechanisms. Unlike the newly proposed sLSTM, these frameworks lack the capacity for memory mixing. This capability of memory mixing is crucial for addressing state tracking tasks, leading LSTMs to possess greater expressive power compared to State Space Models (SSMs) and Transformers~\citep{Merrill:24,Deletang:23}. State tracking is essential for activities such as code evaluation and following entities across extended narratives.

\paragraph{Residually Stacking Architectures.} Similar to the vast majority of modern large-scale deep neural models, xLSTM architectures adopt the approach of residually stacking modular components~\citep{Srivastava:15,He:16}.

This strategy was foundational in the development of deep convolutional models~\citep{He:16} and has been central to the design of Transformers~\citep{Vaswani:17}. The Transformer architecture, in turn, underpins the capabilities of Large Language Models (LLMs) such as GPT-3~\citep{Brown:20short}, ChatGPT~\citep{Schulman:22short}, GPT-4~\citep{Achiam:23gpt4short}, Megatron-LM~\citep{Shoeybi:19}, Gopher~\citep{Rae:21short}, ERNIE 3.0 Titan~\citep{Wang:21arxivshort}, GLaM~\citep{Du:21glamshort}, Chinese M6~\citep{Lin:21}, AlexaTM 20B (multilingual)~\citep{Soltan:22}, OPT~\citep{Zhang:22opt}, Chinchilla~\citep{Hoffmann:22short}, BLOOM~\citep{Scao:22arxivshort}, GLM-130B~\citep{Zeng:22arxivshort}, LaMDA~\citep{Thoppilan:22short}, PaLM~\citep{Chowdhery:22short}, Llama~\citep{Touvron:23arxiv}, and Gemini~\citep{GeminiTeam:23arxiv,Reid:24arxiv}.

\section{Experiments}
\label{sec:Exp}

This section presents a series of empirical analyses of xLSTM, positioning it alongside established approaches in the context of language modeling. Section~\ref{sec:ExpSyntethic} is dedicated to examining the particular strengths of xLSTM using artificial benchmark tasks. Section~\ref{sec:ExpComparison} explores the validation perplexity achieved by several language modeling techniques, each trained with 15B tokens sourced from SlimPajama~\citep{Soboleva:23}, including in-depth ablation studies for xLSTM. We further investigate how these models scale with data size, drawing comparisons to the frameworks described by \citet{Kaplan:20} and \citet{Brown:20short}. 

In Section~\ref{sec:ExpLanguage}, we extend our analysis to a comprehensive language modeling assessment, contrasting xLSTM against top-performing alternatives from Section~\ref{sec:ExpComparison} after their exposure to 300B SlimPajama tokens~\citep{Soboleva:23}. The evaluation encompasses four facets: (1) the ability of models to generalize to extended sequence lengths, (2) validation perplexity and downstream task effectiveness as outlined in~\citep{Sutawika:24}, (3) performance across the 571 domains within the PALOMA language benchmark dataset~\citep{Magnusson:23arxivshort}, and (4) data scaling properties of all models when provided twentyfold greater training volume.

\label{sec:xlstm_blocks_notation}
In every experiment, we denote xLSTM[$a$:$b$] to represent the proportion $a/b$ of mLSTM-based to sLSTM-based xLSTM blocks. As an illustration, xLSTM[7:1] indicates that, among eight total blocks, seven utilize the mLSTM architecture while one employs the sLSTM architecture. In the standard configuration with 48 total blocks, this implies 42 mLSTM-based blocks and 6 sLSTM-based blocks. Additionally, throughout all experiments, we consistently apply up-projection blocks before mLSTM and after sLSTM blocks.

\subsection{Synthetic Tasks and Long Range Arena}
\label{sec:ExpSyntethic}
Initially, we evaluate xLSTM's novel exponential gating paired with memory mixing by applying it to formal language tasks~\citep{Deletang:23}. Next, we investigate how xLSTM's innovative matrix memory contributes to the Multi-Query Associative Recall benchmark~\citep{Arora:23arxiv}. Lastly, we analyze xLSTM's capabilities for handling lengthy input sequences in the Long Range Arena~\citep{Tay:21}.

\paragraph{Evaluation of xLSTM’s Exponential Gating with Memory Mixing.} We investigate the effect of the novel exponential gating and memory mixing mechanisms in xLSTM, which are anticipated to facilitate the resolution of state tracking challenges~\citep{Merrill:24, Merrill:23}. Our approach involves adapting and enhancing the formal language benchmarks from \citet{Deletang:23} by allowing training sequences of varying lengths, thereby supporting length extrapolation. Comprehensive details regarding these tasks and the expanded experimental outcomes are provided in Appendix~\ref{sec:appExpSynthetic-formalLang}. 

For benchmarking, we contrast xLSTM against multiple baselines, such as Transformers, State Space Models, and classical Recurrent Neural Networks, by measuring task-specific token accuracy, rescaled to range from 0 (random guessing) to 1 (error-free performance). We analyze the performance of 2-block configurations, including: xLSTM[0:1] (using only sLSTM blocks), xLSTM[1:0] (utilizing only mLSTM blocks), xLSTM[1:1], Llama, Mamba, RWKV, Retention, Hyena, standard LSTM, and LSTM integrated within Transformer blocks (LSTM (Block)). An overview of the experimental findings is depicted in Figure~\ref{fig:formal-main}.

Architectures lacking explicit memory mixing—such as standard Transformers and State Space Models—demonstrate limitations in addressing certain regular grammar tasks, like the parity benchmark, due to their inability to track state. This observation corroborates prior reports indicating that such models possess fundamentally weaker computational abilities than RNN-based approaches~\citep{Merrill:24, Merrill:23, Deletang:23}.

\paragraph{Test of xLSTM's Memory Capacities on Associative Recall Tasks.} This experiment investigates the matrix memory capabilities of xLSTM using the Multi-Query Associative Recall task~\citep{Arora:23arxiv}, assessing the extent to which the model can retain information. For each input sequence, key-value pairs are sampled at random from a large vocabulary and must be stored by the model for future retrieval. To intensify the challenge beyond the standard task, we increase both the number of key-value pairs up to 256 and the sequence context length up to 2048, thereby enabling a more comprehensive evaluation of each model’s memory strength. We conduct comparisons using 2-block configurations of Llama, Mamba, RWKV-5, RWKV-6, xLSTM[1:1], and xLSTM[1:0], and report the accuracy with which each architecture retrieves stored pairs. Given that Transformers (e.g., Llama) are known to possess a memory scaling exponentially with the coding dimension~\citep{Ramsauer:21}, they serve as a baseline for this benchmark. Experimental outcomes are depicted in Figure~\ref{fig:mqar-main}. Notably, xLSTM[1:1] achieves the highest accuracy among non-Transformer contenders, performing strongly even at lower model sizes. Contrary to possible expectations, integrating an sLSTM block does not reduce memory capability; rather, it appears to augment it, as demonstrated particularly in the most demanding setting with 256 key-value pairs. Further empirical evidence and extrapolation analyses—provided in Appendix~\ref{sec:appExpSynthetic-mqar}—suggest that the superior memory mechanisms of xLSTM also facilitate robust generalization to context lengths exceeding those encountered during model training.

\paragraph{Test of xLSTM's Long Context Capabilities on Long Range Arena.} To evaluate how xLSTM manages extended sequences and sizable contexts, we benchmark multiple models using the Long Range Arena~\citep{Tay:21}. Across every task, xLSTM exhibits robust and reliable results, indicating that its architecture excels at processing and adapting to the diverse challenges presented by long-context scenarios.

For more details, see Appendix~\ref{sec:appExpSynthetic-lra}.

\subsection{Method Comparison and Ablation Study}
\label{sec:ExpComparison}

To explore the primary inquiry posed in this study—namely, the capabilities of our novel LSTM variants in large-scale language modelling—we conduct experiments where xLSTMs, Transformers, State Space Models, and additional approaches are trained on a dataset of 15B tokens sourced from SlimPajama, under uniform auto-regressive conditions. Model performance is evaluated using a shared validation set, and detailed ablation studies are carried out for the xLSTMs.

\paragraph{Assessment of xLSTM Relative to Alternative Approaches.} For comparative analysis, we trained various models on a dataset containing 15 billion tokens sourced from SlimPajama~\citep{Soboleva:23}. Model evaluation was conducted via their perplexity scores obtained on a reserved validation set. The suite of approaches considered includes: xLSTM (the novel technique introduced in this work), GPT-3 (Transformer architecture)~\citep{Brown:20short}, Llama (Transformer)~\citep{Touvron:23arxiv}, H3 (SSM)~\citep{Fu:23}, Mamba (SSM)~\citep{Gu:24arxiv}, RWKV-4 (RNN)~\citep{Peng:23arxivshort}, RWKV-5 (RNN)~\citep{Peng:24arxivshort}, RWKV-6 (RNN)~\citep{Peng:24arxivshort}, GLA (linear Transformer)~\citep{Yang:23arxiv}, HGRN (RNN)~\citep{Qin:23}, HGRN2 (RNN)~\citep{Qin:24arxiv}, RetNet (linear Transformer)~\citep{Sun:23arxiv}, Hyena (linear Transformer)~\citep{Poli:23}, xLSTM[1:0], and xLSTM[7:1] (refer to Section~\ref{sec:xlstm_blocks_notation} for definitions). 

Training was carried out using mixed-precision settings. Notably, RWKV-5, RWKV-6, GLA, and HGRN2 were not trained using PyTorch’s automated mixed-precision approach, as discussed further in Appendix Section~\ref{sec:appGenTrainProc}. The methods are classified as (a) Transformers, (b) State Space Models (SSMs), and (c) Recurrent Neural Networks (RNNs), including linear Transformer variants. Linear Transformers employ alternative mechanisms for attention, distinguishing them from standard Transformer architecture. Model configurations were standardized, mirroring a GPT-3 setup featuring 350M parameters, an embedding dimension of 1024, and 24 residual blocks. It is worth noting that parameter counts are reduced for GPT-3, as it exclusively employs shared weights for token and output embeddings. Table~\ref{tab:spaj15b_model_results} summarizes the findings, showing xLSTM to achieve superior validation perplexity when compared against all other tested approaches. Expanded results and commentary are found in Appendix~\ref{sec:appExpComparison}. Figure~\ref{fig:temp_sclaw15B} illustrates the scaling dynamics for this experimental context, supporting the claim that xLSTM exhibits beneficial scaling characteristics with increasing model size.

\paragraph{Ablation Studies.}
As shown in Table~\ref{tab:spaj15b_model_results} and Figure~\ref{fig:temp_sclaw15B}, xLSTM delivers strong performance in language modeling tasks when trained on 15B tokens sourced from SlimPajama. To systematically assess the contribution of each modification leading from LSTM to xLSTM, we gradually transform a standard LSTM model into the xLSTM architecture. The process begins with incorporating LSTM layers into pre-LayerNorm residual structures. Subsequently, the architecture is enhanced by introducing a post up-projection module. The final stage of this transition includes the application of exponential gating and matrix memory mechanisms.

The results are shown in Table~\ref{tab:ablstudies} (top). 

The ablation analyses indicate that the notable enhancement in performance is the result of the combination of exponential gating and the matrix memory. Moreover, given the critical role of gating in RNNs and State Space Models, we systematically evaluate various gating strategies. The results presented in Table~\ref{tab:ablstudies} (bottom) demonstrate that making each gate trainable and responsive to the input yields incremental benefits. Further experiments exploring the contributions of specific backbone elements are provided in Appendix~\ref{sec:appAblDetails}.

\subsection{xLSTM as Large Language Model}
\label{sec:ExpLanguage}

This study concludes with extensive language modeling experiments aimed at evaluating xLSTM as a large language model. To do so, we scale up the training dataset and utilize 300B tokens from SlimPajama—matching the training regime adopted by models such as Mamba~\citep{Gu:24arxiv} and Griffin~\citep{De:24arxiv}.

xLSTM is benchmarked against RWKV-4, Llama, and Mamba, each representing distinct methodological families as outlined in Section~\ref{sec:ExpComparison}. RWKV-4 is selected to represent RNN-based approaches because appropriate training precision configurations for RWKV-5, RWKV-6, and HGRN2 were determined only after the initiation of the 300B token experiments (refer to Appendix~\ref{sec:appGenTrainProc} for details).

We train models at four different scales: 125M, 350M, 760M, and 1.3B parameters. Each model is evaluated for its ability to extrapolate to longer input sequences and its performance on validation data. Their effectiveness on downstream tasks is measured, and their proficiency in language modeling is tested across 571 text domains included in the PALOMA benchmark. Finally, we examine their scaling laws to analyze performance trends as model size increases.

\paragraph{Sequence Length Extrapolation.}
\label{sec:spaj300B_ctx_extrapolation}
To begin, we evaluate the extrapolation capabilities of large-scale (1.3B parameter) versions of xLSTM, RWKV-4, Llama, and Mamba models in terms of sequence length. Each model is initially trained using a fixed context window of 2048 tokens, after which their performance is assessed with context lengths extending up to 16384 tokens. The results are illustrated in Figure~\ref{fig:spaj300B_seq_extrapolation}. Notably, xLSTM demonstrates a superior ability to preserve low perplexity values as context length increases, distinguishing itself from the other architectures considered.

\paragraph{Validation Perplexity and Downstream Tasks.}
\label{sec:spaj_downstream_eval}
Next, across all scales, we assess xLSTM, RWKV-4, Llama, and Mamba on the SlimPajama validation dataset, targeting both next token prediction performance and several downstream benchmarks that evaluate common sense reasoning abilities.

The third column in Table~\ref{tab:lmeval} presents the perplexity scores on the validation set for various approaches. Notably, xLSTM[1:0] and xLSTM[7:1] demonstrate superior performance across all model sizes when evaluated on this metric. Meanwhile, the remaining columns in Table~\ref{tab:lmeval} report results from downstream task evaluations. xLSTM consistently outperforms other methods for nearly all tasks and model sizes, with the exception of the ARC task, where Mamba occasionally achieves the best results. Additional details can be found in Appendix~\ref{sec:appExpLanguage}.

\paragraph{Performance on PALOMA Language Tasks.} In the third phase of our evaluation, we compare various model architectures—xLSTM, RWKV-4, Llama, and Mamba—by assessing their next token prediction capabilities across all available sizes on PALOMA language tasks~\citep{Magnusson:23arxivshort}. The primary metric used is perplexity, calculated for next token predictions over 571 distinct text domains, spanning sources from nytimes.com to the r/depression subreddit.

Table~\ref{tab:ppleval} presents the perplexity results, with the language modeling categories (first seven columns) separated from more specific domain benchmarks (final five columns). Across these tasks, xLSTM[1:0] consistently surpasses xLSTM[7:1] in predictive accuracy.

Specifically, xLSTM[1:0] outperforms Mamba in 568 out of 571 domains (representing 99.5\%), Llama in 486 domains (85.1\%), and RWKV-4 in 570 domains (99.8\%), based on lower perplexity scores. Detailed breakdowns can be found in Appendix~\ref{sec:appExpLanguage}.

\paragraph{Scaling Laws.} Next, we examine power-law scaling phenomena, which facilitate the prediction of model performance as the parameter count increases~\citep{Kaplan:20,Brown:20short}.
Figure~\ref{fig:temp_sclaw300B} illustrates these scaling trends. While all models demonstrate analogous scaling patterns, their baselines vary. RWKV-4 exhibits the lowest performance, trailed by Llama and Mamba. xLSTM outperforms Mamba, maintaining a performance gap comparable to that between Mamba and Llama. The observed scaling suggests that xLSTM will likely retain its advantage over Transformers and State-Space models as model size expands.

\textbf{Generation Times and Maximal Throughput.} In Figure~\ref{fig:inference_time_generation}, we report the measured generation latency for our xLSTM configurations at the 1.3B parameter setting, while Figure~\ref{fig:inference_time_decoding_throughput} (left) presents their peak throughput. Our evaluation benchmarks these xLSTM variants against comparably-sized versions of Mamba, Llama, and RWKV obtained from HuggingFace, with the Llama model benefiting from a static key-value cache. At the time this evaluation was conducted, it was not feasible to fully compile the Transformer Model cache or to employ \texttt{torch.compile} for Mamba’s compilation. All models involved in the generation time study were run at batch size 1 with a pre-fill of 16, an inference setup that is particularly advantageous for the Transformer architecture. 
The results depicted in Figure~\ref{fig:inference_time_generation} highlight that xLSTM and the recurrent baselines, Mamba and RWKV-4, display linear increases in generation time as the sequence length grows, in contrast with the Llama model’s quadratic scaling. In terms of throughput, Figure~\ref{fig:inference_time_decoding_throughput} (right) reveals that xLSTM can accommodate significantly larger batch sizes than Llama, owing to its constant memory requirements, which leads to the greatest observed throughput among the models considered. For this analysis, various batch sizes and pre-fills are explored, particularly for the Llama benchmark.

\section{Limitations}

(i) Unlike mLSTM, the sLSTM’s mechanism for mixing memory inherently restricts parallel computation, making a fast parallel implementation infeasible. Despite this, we created an efficient CUDA kernel for sLSTM, which currently operates at under twice the runtime of our parallelized mLSTM version. (ii) The existing CUDA kernels for mLSTM lack thorough optimization, rendering them approximately four times slower compared to FlashAttention and the scan method utilized in Mamba. With more optimized CUDA kernels, similar to those developed for FlashAttention, substantial speed gains could be achieved. (iii) Since mLSTM relies on processing $d \times d$ matrix memory, the computational demand is significant. However, updates and retrievals from the memory do not require additional parameters and benefit from the parallelizability afforded by conventional matrix operations, thereby making the increased runtime marginal relative to the complexity of the memory structure. (iv) Proper selection of the forget gate initialization parameters is crucial. (v) As the matrix memory in mLSTM is invariant to sequence length, extending sequence lengths can potentially overwhelm the memory when contexts grow larger. Nevertheless, empirical results indicate this issue does not hinder performance with contexts up to 16k tokens (see Section~\ref{sec:spaj300B_ctx_extrapolation}). (vi) Given the high computational cost of large language model experiments, neither the architecture nor hyperparameters were fully refined, particularly in the context of larger xLSTM variants. We expect substantial improvements in xLSTM performance will require rigorous architectural and hyperparameter optimization.

\section{Conclusion}
Our investigation has provided a partial response to the central question: How effective is language modeling when scaling LSTM architectures to billions of parameters? Thus far, our results suggest that such scaling enables performance comparable to that of leading approaches such as Transformers and State Space Models. We introduced xLSTM, which incorporates exponential gating combined with memory mixing and an innovative memory architecture. Across language modeling tasks, xLSTM demonstrates strong results relative to state-of-the-art techniques, including Transformers and State Space Models. Furthermore, analysis of scaling behavior suggests that even larger xLSTM architectures could challenge the dominance of existing Transformer-based Large Language Models. The versatility of xLSTM may also bring significant advances to fields including Reinforcement Learning, Time Series Prediction, and the modeling of physical systems.

\end{document}